% GENERATED FILE. Edit canonical lessons, then run npm run generate:books.
% canonical-source-hash: fnv1a64:e60917dfdfce3ebd
% canonical-lessons: MR-C507-bharti, MR-C507-rojgar, MR-C507-prasiksan, MR-C507-hatoti, MR-C507-ruju-hone

\chapter{Finding Work}
\label{ch:mr-finding-work}
\hlchaptermodality{507}

\begin{chapteropening}
\textbf{By the end of this chapter:} \emph{I can talk about looking for work, training and starting a job.}

Complete the last lesson of chapter 507: I can talk about looking for work, training and starting a job.
\end{chapteropening}

\section[\texorpdfstring{\mr{भरती} (bhartī)}{भरती (bhartī)}]{\mr{भरती} (bhartī) — recruitment; enrolment}

\label{lesson:MR-C507-bharti}



\begin{quote}
Before the new one: say the Marathi for a document, then the Marathi for to scrutinize.
\end{quote}

\subsection*{You'll want to know: \mr{भरती}}
\textbf{\mr{भरती}} — \emph{bhartī} — \textquotedblleft{}recruitment; enrolment\textquotedblright{}.

\begin{quote}
\mr{पोलिसांची} \mr{भरती} \mr{पुढच्या} \mr{महिन्यात} \mr{आहे}.
\end{quote}

\begin{quote}
\emph{polisāñcī bhartī puḍhcyā mahinyāt āhe.}
\end{quote}

\begin{quote}
The police recruitment is next month.
\end{quote}

Feminine.

\mr{भरती} \mr{होणे}, to be recruited; \mr{भरती} \mr{करणे}, to recruit.

\begin{cousinweb}
\textbf{\mr{भरती}} (\emph{bhartī}) is built on \mr{भरणे}, \textquotedblleft{}to fill\textquotedblright{}: a recruitment fills the empty places. The verb is inherited from Sanskrit \emph{bharati}, \textquotedblleft{}bears, fills\textquotedblright{} (CDIAL 9397).
\end{cousinweb}

\begin{grammarlens}[title={What you've built}]
One more word for describing what you lived and felt.
\end{grammarlens}

\subsection*{Your turn}
\begin{itemize}
\raggedright
  \item \emph{Say it:} \emph{bhartī}
  \item \emph{Say it:} \emph{bhartī}, once more
  \item \emph{Say it:} \emph{polisāñcī bhartī puḍhcyā mahinyāt āhe.} — The police recruitment is next month.
  \item \emph{Recall:} say the Marathi for a document, then the Marathi for to scrutinize, then say \emph{bhartī} again
\end{itemize}

\begin{tcolorbox}[breakable,colback=teal!4,colframe=teal!35!black,title={Before you move on}]
What does \mr{भरती} mean? (\textquotedblleft{}Recruitment; enrolment\textquotedblright{}.) Say it once more.
\end{tcolorbox}
\section[\texorpdfstring{\mr{रोजगार} (rojgār)}{रोजगार (rojgār)}]{\mr{रोजगार} (rojgār) — employment; a livelihood}

\label{lesson:MR-C507-rojgar}



\begin{quote}
Before the new one: say the Marathi for to scrutinize, then the Marathi for recruitment.
\end{quote}

\subsection*{You'll want to know: \mr{रोजगार}}
\textbf{\mr{रोजगार}} — \emph{rojgār} — \textquotedblleft{}employment; a livelihood\textquotedblright{}.

\begin{quote}
\mr{शहरात} \mr{रोजगार} \mr{शोधायला} \mr{अनेक} \mr{तरुण} \mr{येतात}.
\end{quote}

\begin{quote}
\emph{śahrāt rojgār śodhāylā anek taruṇ yetāt.}
\end{quote}

\begin{quote}
Many young people come to the city to look for work.
\end{quote}

Masculine.

\begin{cousinweb}
\textbf{\mr{रोजगार}} (\emph{rojgār}) is a loan from Persian \emph{rōzgār}, \textquotedblleft{}time, fortune, livelihood\textquotedblright{}, built on \emph{rōz}, \textquotedblleft{}day\textquotedblright{}.
\end{cousinweb}

\begin{grammarlens}[title={What you've built}]
One more word for describing what you lived and felt.
\end{grammarlens}

\subsection*{Your turn}
\begin{itemize}
\raggedright
  \item \emph{Say it:} \emph{rojgār}
  \item \emph{Say it:} \emph{rojgār}, once more
  \item \emph{Say it:} \emph{śahrāt rojgār śodhāylā anek taruṇ yetāt.} — Many young people come to the city to look for work.
  \item \emph{Recall:} say the Marathi for to scrutinize, then the Marathi for recruitment, then say \emph{rojgār} again
\end{itemize}

\begin{tcolorbox}[breakable,colback=teal!4,colframe=teal!35!black,title={Before you move on}]
What does \mr{रोजगार} mean? (\textquotedblleft{}Employment; a livelihood\textquotedblright{}.) Say it once more.
\end{tcolorbox}
\section[\texorpdfstring{\mr{प्रशिक्षण} (praśikṣaṇ)}{प्रशिक्षण (praśikṣaṇ)}]{\mr{प्रशिक्षण} (praśikṣaṇ) — training}

\label{lesson:MR-C507-prasiksan}



\begin{quote}
Before the new one: say the Marathi for recruitment, then the Marathi for employment.
\end{quote}

\subsection*{You'll want to know: \mr{प्रशिक्षण}}
\textbf{\mr{प्रशिक्षण}} — \emph{praśikṣaṇ} — \textquotedblleft{}training\textquotedblright{}.

\begin{quote}
\mr{नोकरीच्या} \mr{आधी} \mr{मला} \mr{एक} \mr{महिना} \mr{प्रशिक्षण} \mr{मिळालं}.
\end{quote}

\begin{quote}
\emph{nokrīcyā ādhī malā ek mahinā praśikṣaṇ miḷālaṁ.}
\end{quote}

\begin{quote}
Before the job I got a month of training.
\end{quote}

Neuter.

The trainer is \mr{प्रशिक्षक}.

\begin{cousinweb}
\textbf{\mr{प्रशिक्षण}} (\emph{praśikṣaṇ}) is a tatsama coinage from Sanskrit \emph{pra-}, \textquotedblleft{}forth\textquotedblright{}, and \emph{śikṣaṇa}, \textquotedblleft{}teaching, instruction\textquotedblright{} (Monier-Williams).
\end{cousinweb}

\begin{grammarlens}[title={What you've built}]
One more word for describing what you lived and felt.
\end{grammarlens}

\subsection*{Your turn}
\begin{itemize}
\raggedright
  \item \emph{Say it:} \emph{praśikṣaṇ}
  \item \emph{Say it:} \emph{praśikṣaṇ}, once more
  \item \emph{Say it:} \emph{nokrīcyā ādhī malā ek mahinā praśikṣaṇ miḷālaṁ.} — Before the job I got a month of training.
  \item \emph{Recall:} say the Marathi for recruitment, then the Marathi for employment, then say \emph{praśikṣaṇ} again
\end{itemize}

\begin{tcolorbox}[breakable,colback=teal!4,colframe=teal!35!black,title={Before you move on}]
What does \mr{प्रशिक्षण} mean? (\textquotedblleft{}Training\textquotedblright{}.) Say it once more.
\end{tcolorbox}
\section[\texorpdfstring{\mr{हातोटी} (hātoṭī)}{हातोटी (hātoṭī)}]{\mr{हातोटी} (hātoṭī) — a knack; a deft touch}

\label{lesson:MR-C507-hatoti}



\begin{quote}
Before the new one: say the Marathi for employment, then the Marathi for training.
\end{quote}

\subsection*{You'll want to know: \mr{हातोटी}}
\textbf{\mr{हातोटी}} — \emph{hātoṭī} — \textquotedblleft{}a knack; a deft touch\textquotedblright{}.

\begin{quote}
\mr{माझ्या} \mr{आजीला} \mr{भाकरी} \mr{करण्याची} \mr{चांगली} \mr{हातोटी} \mr{आहे}.
\end{quote}

\begin{quote}
\emph{mājhyā ājīlā bhākrī karṇyācī cāṅglī hātoṭī āhe.}
\end{quote}

\begin{quote}
My grandmother has a real knack for making bhakri.
\end{quote}

Feminine.

The one who has it takes -\mr{ला}: \mr{तिला} \mr{हातोटी} \mr{आहे}.

\begin{cousinweb}
\textbf{\mr{हातोटी}} (\emph{hātoṭī}) is built on \mr{हात}, \textquotedblleft{}hand\textquotedblright{}, inherited from Sanskrit \emph{hasta} (CDIAL 14024): a knack sits in the hands.
\end{cousinweb}

\begin{grammarlens}[title={What you've built}]
One more word for describing what you lived and felt.
\end{grammarlens}

\subsection*{Your turn}
\begin{itemize}
\raggedright
  \item \emph{Say it:} \emph{hātoṭī}
  \item \emph{Say it:} \emph{hātoṭī}, once more
  \item \emph{Say it:} \emph{mājhyā ājīlā bhākrī karṇyācī cāṅglī hātoṭī āhe.} — My grandmother has a real knack for making bhakri.
  \item \emph{Recall:} say the Marathi for employment, then the Marathi for training, then say \emph{hātoṭī} again
\end{itemize}

\begin{tcolorbox}[breakable,colback=teal!4,colframe=teal!35!black,title={Before you move on}]
What does \mr{हातोटी} mean? (\textquotedblleft{}A knack; a deft touch\textquotedblright{}.) Say it once more.
\end{tcolorbox}
\section[\texorpdfstring{\mr{रुजू} \mr{होणे} (rujū hoṇe)}{रुजू होणे (rujū hoṇe)}]{\mr{रुजू} \mr{होणे} (rujū hoṇe) — to join a post; to report for duty}

\label{lesson:MR-C507-ruju-hone}



\begin{quote}
Before the new one: say the Marathi for training, then the Marathi for a knack.
\end{quote}

\subsection*{You'll want to know: \mr{रुजू} \mr{होणे}}
\textbf{\mr{रुजू} \mr{होणे}} — \emph{rujū hoṇe} — \textquotedblleft{}to join a post; to report for duty\textquotedblright{}.

\begin{quote}
\mr{अरुण} \mr{सोमवारी} \mr{नवीन} \mr{कार्यालयात} \mr{रुजू} \mr{झाला}.
\end{quote}

\begin{quote}
\emph{aruṇ somvārī navīn kāryālayāt rujū jhālā.}
\end{quote}

\begin{quote}
Arun joined the new office on Monday.
\end{quote}

Past (he …): \mr{तो} \mr{कामावर} \mr{रुजू} \mr{झाला}, \emph{to kāmāvar rujū jhālā}.

The place takes -\mr{त} or -\mr{वर}: \mr{कार्यालयात} \mr{रुजू} \mr{होणे}, \mr{कामावर} \mr{रुजू} \mr{होणे}.

\begin{cousinweb}
\textbf{\mr{रुजू} \mr{होणे}} (\emph{rujū hoṇe}) is \mr{रुजू}, a loan usually traced to Arabic \emph{rujūʿ}, \textquotedblleft{}turning to\textquotedblright{}, through Persian, with \mr{होणे}: to present oneself for duty.
\end{cousinweb}

\begin{grammarlens}[title={What you've built}]
One more word for describing what you lived and felt.
\end{grammarlens}

\subsection*{Your turn}
\begin{itemize}
\raggedright
  \item \emph{Say it:} \emph{rujū hoṇe}
  \item \emph{Say it:} \emph{rujū hoṇe}, once more
  \item \emph{Say it:} \emph{aruṇ somvārī navīn kāryālayāt rujū jhālā.} — Arun joined the new office on Monday.
  \item \emph{Recall:} say the Marathi for training, then the Marathi for a knack, then say \emph{rujū hoṇe} again
\end{itemize}

\begin{tcolorbox}[breakable,colback=teal!4,colframe=teal!35!black,title={Before you move on}]
What does \mr{रुजू} \mr{होणे} mean? (\textquotedblleft{}To join a post; to report for duty\textquotedblright{}.) Say it once more.
\end{tcolorbox}
